\section{Restatement of Purpose}

To whom it may concern.

Someone suggested there needs to be some ``dialog", presumably between Gornahoor and some unidentified dialogist. That presumes there is something to discuss, necessitating some common ground, and resulting in a common goal. The ground is Tradition. The goal is the recovery of Tradition in the West. We have not deviated from that in over five years.

In January 2008, we set the stage in one of the first posts here on the Prophecies of \textbf{Rene Guenon}, Part 1 and Part 2. Some commenters seem to believe we are trying to be persuasive although we have been clear that we are only interested in facts, their consequences, and the various possibilities in give situation. These are the two relevant facts.

\begin{itemize}
\item The West had a valid Tradition up until the Medieval Era 
\item The Tradition was eventually lost, after the destruction of the Templars according to Guenon, and perhaps a little later in Evola's view 
\end{itemize}
The following corollary immediately follows from those two observations:

\begin{itemize}
\item In order to fully understand Tradition, it is sufficient to look to medieval Europe, its spiritual tradition, and its social structure 
\end{itemize}
In other words, it is not necessary to look to alien traditions, but rather to Europe's intellectual and spiritual past. This was brought to our attention by an essay by \textbf{Ananda Coomaraswamy}, in which he made clear that a European could not understand the Vedanta unless he was able to understand the great traditional writers of the West. There is a corollary to that observation:

\begin{itemize}
\item Anyone who scoffs at the spiritual tradition of medieval Europe is unlikely to recognize or understand the traditional elements in the Vedanta. 
\end{itemize}
A project immediately presents itself, one which Gornahoor has endeavored to fulfill:

\begin{itemize}
\item Investigate the medieval spiritual tradition and demonstrate its teachings on metaphysics, psychology, anthropology, history, sociology, politics, and so on. 
\end{itemize}
Now, all this is more than a mental exercise. The ultimate purpose is to point out a possible way for the recovery of Tradition in the West. Moreover, that purpose is not the result of a merely arbitrary and subjective preference. As the modern world cannot continue to sustain itself on its present path, Tradition will necessarily return. However, that return can be relatively peaceful or quite painful. As noted in the links to Guenon's predictions, there are three possibilities:

\begin{itemize}
\item Degeneration. If no effort is made to restore Tradition, the West will degenerate into barbarity. \textbf{Carl Jung} noticed that in the psychology of his patients. Concerning the abandonment of the Western spiritual tradition by modern men, he had this to say:

\begin{quotex}
At a time when a large part of mankind is beginning to discard Christianity, it may be worth our while to try to understand why it was accepted in the first place. It was accepted as a means of escape from the brutality and unconsciousness of the ancient world. As soon as we discard it, the old brutality returns in force … This is not a step forwards, but a long step backwards into the past.
\end{quotex}
\item Assimilation. Outsiders from the East will bring Tradition back to Europe, whether by consent or by force. In the former case, Europeans would adopt the Traditions of the newcomers. In the latter, the Tradition would come by force marked by painful ethnic or racial conflicts. 
\item Transformation. In this scenario, the West returns to Tradition spontaneously and voluntarily. This would be initiated by a group of men that Guenon calls the ``elite", in the sense that they understand tradition and ``true intellectuality". They would be the initial leavening that would eventually spread. 
\end{itemize}
Since we are not interested in promoting the cause of degeneration, and since we are not ``outsiders", Gornahoor's real project is almost too obvious to mention.

\begin{itemize}
\item To promote the spontaneous and voluntary transformation of the West. 
\end{itemize}
There are several preparatory steps required. Here are some we are engaged in:

\begin{itemize}
\item Make clear what the elite must know, both intellectually and morally. That is why we have provided so many translations into English of rare works and have brought to light many ancient texts. Obviously, we read them in the light of Tradition and not as moderns who have little conception of their real meaning. 
\item Encourage readers to engage with writings that are known to be initiatory, such as Dante, The Romance of the Rose, Boccaccio, etc. 
\item Investigate spiritual and intellectual currents that still retain real and recoverable traces of tradition. 
\end{itemize}
There are other initiatives in motion, but it is too early to discuss them.

There is absolutely nothing new or novel in all this, it all derives from Guenon and other writers on Tradition. Hence, there is really nothing to dialog about; you either accept the premise or you don't. We are interested only in the former. It is not dialogists that are required, but rather those who want to develop such ideas further, based on their own abilities and interests.



\flrightit{Posted on 2013-10-14 by Cologero }

\begin{center}* * *\end{center}

\begin{footnotesize}\begin{sffamily}



\texttt{Thorsten on 2013-10-14 at 10:34 said: }

Cologero,

thanks for the clarification. When I suggested a dialogue I was referring primarily to the fact that self-styled ``traditionalists" in the New Right/Alt Right need to wiling to be more open to traditional doctrine, especially in the shape it assumed in Medieval Europe, and to be more open to metaphysical reality in general. (Otherwise they really have no business calling themselves traditionalists.)

On the other hand I would like to humbly suggest that as prideful and sensually/aesthetically oriented [What else is left in the absence of a living tradition than tribal and racial identity and old symbolic language (e.g. art) from a bygone civilization? These seem to be the default surrogates for tradition, and the only remaining alternatives to chaos when you were raised outside of tradition or have been heretofore fed such a degraded version of tradition that you cannot in good conscience submit to it.] these apostates are that they are not all beyond saving as a group, and should be handled as misguided rather than willfully contrarian. Now doubtless a number of the most strident voices in these circles are in fact willfully contrarian, puerile, and deliberately disrespectful of Western tradition to a degree that is shameful and they are probably not worth the time of Gornahoor or any other advocates of western tradition.

I'm confident that most of the New Right/Alt Right are not interested in Muslim overlordship and therefore are opposed the the option of Assimilation, although I'm not so sure that they would all be opposed to some level brutality and barbarity in terms of a more militarized society and harsher corporal punishment (It's undeniable that western society was markedly more brutal and barbarous until quite recently in this regard.) If brutality and barbarism refers to coarseness of taste, vulgarity, and surrender to hedonism I think its obvious that most of them would oppose Degeneration in this sense, as this is frequently the subject of articles of the New Right. On the other hand I think the third possibility of Transformation might actually be received quite well as this would appeal to those who were not born into an existing tradition or have been alienated the contemporary Chrisitianity of any sort because they were raised on some anti-traditional low church protestantism which is devoid traditional values such as inequality, metaphysics as opposed to literalist readings of scripture, and the primacy of the intellect as opposed to blind faith.

I myself am not a member in any sense of the New Right, but I do think that there is more potential among those who frequent such circles than there appears to be if you only focus on the stance of the writers and editors. As a former Nietzsche-maniac, and stubborn materialist like so many of those frequent and comment on such websites I think that many of these iconoclasts could be esotericists in the making.

Kind regards,

Thorsten


\hfill

\texttt{Thomas Blanchard on 2013-10-14 at 14:01 said: }

Thorsten,

I'll chime in since I am loosely involved in various New Right spheres, and can echo most of your feelings on this, although I think I've come to a different understanding of things.

I agree that the reformation of an authentic new second estate (warrior caste) will be a necessary step in an overall return to order in the west, just as it will be necessary to reform the third estate on healthy principles. Gornahoor's primary mission (and someone correct me if I'm wrong) is the recovery of the western tradition and and the spiritual development of individuals who will eventually form a truly legitimate first estate that will serve as its elite vanguard.

This being said, a parallel effort to build a true second estate does make sense, but they key is that it should be properly subordinated to a new first estate; the castes and vocations should have proper hierarchical relations with one another. All of the various New Right movements of today represent different manifestations of this impulse, but without a proper spiritual authority to guide them (and you are correct that the modern exoteric Church as the world sees it is rather weak and uninspiring), these efforts constitute a rebellious attempt to re-establish the Bronze Age (written of by Hesiod). In this case I respectfully disagree with Cologero's non-distinction between pure modernists and the New Right types that one might find at a place like Counter-Currents: modernists represent the Iron Age, the ``Renaissance", the French Revolution, and the ``Enlightenment", but the New Right truly wishes for the return of the Bronze Age and their writers are men like Homer and Nietzsche.

In a sense, the dialogue you have been talking about has been taking place for a long time – you can see it in Plato's works: Socrates' debates with Thucydides or especially with Callicles, where his opponents represent these more ``New Right" sorts.

As a person who also has mixed proclivities and tends to sympathize with Nietzsche, I share your feelings, but it's best for the time being to focus on re-establishing a spiritual elite before trying to give direction to the chaotic mess of political agitators who are united only in their opposition to the Iron Age.

For your personal development, if you cannot commit yourself to a religion now, at least I would recommend taking up things like traditional martial arts and mountain climbing – Evola wrote primarily for people of our type, and some of his last advice to young people was not to worry about forming an ``order", but to simply focus on building character. I believe that this is good advice.


\hfill

\texttt{Thorsten on 2013-10-14 at 20:06 said: }

Thomas,

I'm on the same page with you and Gornahoor on the primacy of establishing a reformed spiritual/intellectual elite/oratores/brahmin caste (call it what you will). Thought precedes action and the wise must rule the strong and I say that is all the more reason to quash counterfeit spiritualities propagated by some in the new right and to formulate a working alternative. How to do this specifically I am not sure and won't presume to tell anyone.

I myself cannot say with certainty whether I am more a bellator than an orator but I do have some experience with rock climbing and I intend to take up some form of martial arts when I have the means to do so. Anyway I understand the necessity of rebuilding from the top down as has been reiterated here at Gornahoor many times, but if we don't get things rebuilt (bring about the Transformation) before society collapses into utter degeneracy do we simply allow some semblance of order to be imposed on the West by Wahhabi radicals (Assimilation) by means of brutality and barbarism which may make the barbarian incursions in the last days of the Roman empire look merciful in comparison?

To be clear am no enemy of Islam. I am only pointing out that when the present ideology of materialism and egalitarianism implodes upon itself the most consistent (and forceful) ideology in the West will not only be Muslim but the most ferocious variant thereof.


\hfill

\texttt{Michael on 2013-10-14 at 20:25 said: }

Cologero, I understand the concept of an elite, but I don't understand how the elite will act as a ``leaven" to transform society. Is it purely through contemplation? Or will there also be action similar to the way that enlightenment ideas spread through society? From reading Guenon and Tomberg, I have the impression that it is an unseen spiritual leavening process.


\hfill

\texttt{Cologero on 2013-10-14 at 23:18 said: }

@ Thomas Blanchard Those are mostly excellent suggestions … more generally, besides rock climbing, there are other elements to master depending on your personal equation. See Qualifications for Initiation for example. If only the NR would follow your advice and stay out of religious and spiritual matters. Their proper domain is the manifestation which should be the focus. An important point you bring up is the ``third estate", which does not get the attention it needs. The modern world is in the hands of that estate, not some decadent spiritual tradition as the NR seems to be obsessed with.

There is some subtleties with Socrates, who was executed for atheism and impiety. His call for the rule of philosophers in touch with the ideas replaced the Homeric vision of the rule of the gods. So you could see him as the precursor of modernity on the one hand, or else as someone who spread esoteric teachings prematurely.


\hfill

\texttt{Cologero on 2013-10-14 at 23:32 said: }

@Thorsten

I put up the quotes from Alice Bailey to show how events happen on their own accord apart from the wills of men. She saw turmoil inaugurating a new age of one religion, one race, etc, although it is hard to be so sanguine about that. It is odd that the Theosophists seen eternal progress instead of the cycles of decline, considering their reliance on Asian teachings.

Back to the point, events will unfold as they do. Decisions made and actions executed today will have consequences down the road. I don't know how Guenon was able to see the possibly imminent ethnic conflict in 1921, but they are certainly foreseeable now. Yet it doesn't have to be that way.

I don't know why you feel any personal agitation or anxiety. It is not as though you are in charge of anything nor is the worst likely to happen in your lifetime. Besides rock climbing, seek out Seneca's 41st letter to Lucillius:

\begin{quotex}
If you see a man who is unterrified in the midst of dangers, untouched by desires, happy in adversity, peaceful amid the storm, who looks down upon men from a higher plane and view the gods on a footing of equality, will not a feeling of reverence for him steal over you? …

When a soul rises superior to other souls, when it is under control, when it passes through every experience as if it were of small account, when it smiles at our fears and at our prayers, it is stirred by a force from heaven. A thing like this cannot stand upright unless it be propped by the divine. 

\end{quotex}
That is the ideal to strive for. Jung's commentary on that passage is also quite interesting, but I'll save it for another time.


\hfill

\texttt{Cologero on 2013-10-14 at 23:44 said: }

Guenon ties it to a very specific event: the destruction of the Templars. The reason for this may be difficult to understand or even to accept, but it has to do both with the exodus of esoteric orders from Europe and the ties between esoteric Christian and esoteric Islamic orders. Guenon has left a very interesting clue about that, although I am not ready to make it public without understanding it more deeply. I have pointed out the clue on Gornahoor although it will take some digging.

I don't believe Evola dates it so precisely since he did not have the same interest in that topic as Guenon. Instead he refers to the regression of the castes over a longer time. As we have often pointed out, Evola said he believed what every well-bred person prior to the French Revolution considered normal and healthy.


\hfill

\texttt{Cologero on 2013-10-14 at 23:54 said: }

I don't understand either, Michael, as the Spirit will blow where he wills. Yes, false ideas spread quite easily, but true ideas are longer lasting. There is always the feeling of having to ``do something" but instead we need that interior silence, not unlike what Seneca described, which connects us to Heaven which allows the influx of forces.


\hfill


\end{sffamily}\end{footnotesize}
